% Requires booktabs, array, longtable.
{\small
\begin{longtable}{>{\raggedright\arraybackslash}p{68.49pt}>{\raggedright\arraybackslash}p{91.31pt}>{\raggedright\arraybackslash}p{114.14pt}>{\raggedright\arraybackslash}p{182.63pt}}
\caption{Conditional parameter-interval coverage}\label{tab:table_s10}\\
\toprule
\textbf{Months} & \textbf{Parameter} & \textbf{Covered /\allowbreak{} histories} & \textbf{Coverage (\%) [95\% MC CI]} \\
\midrule\endfirsthead
\toprule
\textbf{Months} & \textbf{Parameter} & \textbf{Covered /\allowbreak{} histories} & \textbf{Coverage (\%) [95\% MC CI]} \\
\midrule\endhead
\bottomrule\endfoot
12 & xi & 182 /\allowbreak{} 200 & 91.00 [86.15, 94.58] \\
12 & sigma & 190 /\allowbreak{} 200 & 95.00 [91.00, 97.58] \\
24 & xi & 185 /\allowbreak{} 200 & 92.50 [87.93, 95.74] \\
24 & sigma & 182 /\allowbreak{} 200 & 91.00 [86.15, 94.58] \\
60 & xi & 190 /\allowbreak{} 200 & 95.00 [91.00, 97.58] \\
60 & sigma & 191 /\allowbreak{} 200 & 95.50 [91.63, 97.92] \\
\end{longtable}
\noindent{\footnotesize Separate inherited recovery experiment, 200 histories per setting; not the primary experiment. Coverage conditions on the selected threshold. Clopper--Pearson intervals describe Monte Carlo uncertainty in coverage.\par}
}
